%% kapa: exam=1 round=20 year=2009 period=1 form=A subject=civil_law subject_ko=민법 seq=2 number=2
%% source: Q-Net 감정평가사 기출문제 (cst00309) artlSeq=5024178 file=extracted/1008/5024178/2009년 제20회 감정평가사 1차시험 기출문제/1교시.hwp
%% conversion: kordoc 4.24.51 -> tools/gen1.py (auto); review=proofread-sonnet; flags=none
\begin{question}{2}
물권적 청구권에 관한 설명으로 옳지 \uline{않은} 것은?
\begin{choices}
  \item 상대방의 고의ㆍ과실이 없더라도 물권적 청구권을 행사할 수 있다.
  \item 물권적 청구권은 물권과 분리하여 이를 양도할 수 없다.
  \item 물권적 청구권과 불법행위로 인한 손해배상청구권은 병존할 수 있다.
  \item 물권적 청구권이 물권과 별도로 소멸시효에 걸리는지에 관하여는 학설이 대립한다.
  \item 민법은 소유권에 기한 물권적 청구권을 질권에 준용하는 규정을 두고 있다.
\end{choices}
\end{question}
\answer{5}
